\documentclass[11pt,letterpaper]{article}
\usepackage{fontspec}
\usepackage{tocloft}
\usepackage{multicol}
\usepackage{fancyhdr}
\usepackage{nicefrac}
\usepackage[left=1.3in,right=1.3in,top=1in,bottom=1in]{geometry}
\setmainfont[Scale=1.3,AutoFakeBold=1.5,AutoFakeSlant=0.3]{Doves Type}

\widowpenalty=10000
\clubpenalty=10000
\interlinepenalty=500

\title{Green Chile Chicken Tinga Enchiladas}
\author{}
\date{}

% Configure fancy header
\pagestyle{fancy}
\fancyhf{} % Clear all header and footer fields
\fancyhead[C]{\textit{Green Chile Chicken Tinga Enchiladas}} % Center header with recipe title
\fancyfoot[C]{\thepage} % Center page number at bottom
\renewcommand{\headrulewidth}{0pt} % Remove header line
\renewcommand{\footrulewidth}{0pt} % Remove footer line

\begin{document}

\maketitle
\thispagestyle{empty}

% \begin{quote}\textit{Air-fried Anaheim, poblano, and jalapeño chiles are blended skins-on with canned tomatillos, onion, and garlic, then simmered into a green sauce. Two cups become a creamy Monterey Jack sauce around chicken tinga, charred corn, and raw onion, finished with crumbled cotija, and about 3~cups of plain sauce are kept for later. Makes 12 enchiladas and serves 4--6.}\end{quote}

\section*{Ingredients}
\setlength{\columnsep}{20pt}
\begin{multicols}{2}
\noindent
    Anaheim chiles \dotfill 18~oz. \\
    Poblano chiles \dotfill 6~oz. \\
    Jalapeño \dotfill 1 \\
    Tomatillos, canned \dotfill 1 (28~oz.) can \\
    White onions, medium (divided) \dotfill 2 \\
    Garlic cloves \dotfill 4 \\
    Vegetable oil, divided \dotfill \nicefrac{1}{2}~cup \\
    Ground cumin \dotfill 1\nicefrac{1}{2}~tsp. \\
    Mexican oregano \dotfill 1~tsp. \\
    Kosher salt \dotfill to taste \\
    Lime \dotfill 1 \\
    \columnbreak
    Salted butter \dotfill 3~Tbsp. \\
    All-purpose flour \dotfill 3~Tbsp. \\
    Chicken broth \dotfill 1~cup \\
    Mexican crema \dotfill \nicefrac{1}{2}~cup \\
    Monterey Jack cheese, shredded \dotfill 10~oz. \\
    Cotija cheese \dotfill 3~oz. \\
    Chicken tinga \dotfill 3~cups \\
    Corn kernels \dotfill 1~cup \\
    Fresh cilantro, chopped \dotfill \nicefrac{1}{4}~cup \\
    Corn tortillas (6-inch) \dotfill 12 \\
\end{multicols}

\filbreak
\section*{Charred Chiles}

\noindent
Heat the air fryer to \textit{400°F}

\begin{enumerate}
    \item Rub the \textbf{Anaheim~chiles}, the \textbf{poblano~chiles}, and the \textbf{jalapeño} with about 1~tsp. of the \textbf{vegetable~oil}. Cook them in a single layer at \textit{400°F} for \textit{15~minutes}, turning at \textit{7~minutes}. They are done when the skins are blistered and blackened in patches and the flesh collapses under a gentle press. Pull the \textbf{jalapeño} as soon as it is soft. Continue in \textit{2--3~minute} passes for any chile that is still firm.

    \item Cook a second batch the same way for any chiles that did not fit in the first layer.

    \item When the chiles are cool enough to handle, slit each one and scrape out the seeds. Leave the charred skins on. Add them to the blender.
\end{enumerate}

\filbreak
\section*{Green Sauce}

\noindent
While the chiles cook ---
Drain the \textbf{tomatillos} and discard the liquid; put the fruit in the blender ---
Dice the \textbf{white~onions}; put \nicefrac{1}{2}~cup in \textit{Large Bowl~\#2} (filling) and the rest in the blender ---
Peel the \textbf{garlic} and add it to the blender ---
Combine 1\nicefrac{1}{2}~tsp. \textbf{cumin} and 1~tsp. \textbf{Mexican~oregano} in \textit{Small Bowl~\#1} (sauce spices) ---
Juice the \textbf{lime} into \textit{Small Bowl~\#2}

\begin{enumerate}
    \item Blend the chiles, \textbf{tomatillos}, \textbf{white~onion}, and \textbf{garlic} until the puree is smooth and no papery skin remains, about \textit{45--90~seconds}. Scrape the blender and blend for another \textit{20--30~seconds} if flakes remain. Add a spoonful of water only if the blades stall.

    \item Set a large saucepan over medium heat. Add 2~Tbsp. \textbf{vegetable~oil} and heat until it shimmers. Stir in the \textbf{cumin} and \textbf{Mexican~oregano} (\textit{Small Bowl~\#1}) for \textit{30~seconds}, until the raw spice smell turns fragrant. Pour in the puree. It will sputter. Stir until it settles, and scrape the bottom clean.

    \item Leave the pot at a steady medium simmer. Char the \textbf{corn} in the next step, then stir the sauce again. Keep simmering, stirring often, for \textit{15--25~minutes} in all, until the color has darkened, the raw onion edge is gone, and a finger path through the sauce on a spoon stays open. The surface looks glossy, and the spoon still moves through the pot. Continue for \textit{3--5~minutes} if it still tastes sharp or runs thin.

    \item While the sauce simmers, pat the \textbf{corn} dry. Heat a 12-inch skillet over medium-high heat. Cook the \textbf{corn} in the dry skillet, stirring, for \textit{6--10~minutes}, until the moisture cooks off, the pan looks dry, and the kernels are spotted deep brown and smell toasted. Continue for \textit{2~minutes} if the pan is still steaming or the kernels are pale. Add the \textbf{corn} to \textit{Large Bowl~\#2}.

    \item Take the sauce off the heat. Stir in the \textbf{lime~juice} (\textit{Small Bowl~\#2}) and 1\nicefrac{1}{2}~tsp. \textbf{kosher~salt}. Taste, and add more \textbf{kosher~salt} if it tastes flat. Scrape the sauce into \textit{Large Bowl~\#1}. There should be about 5~cups. Measure 2~cups for the creamy sauce and leave the rest in the bowl. That remainder, about 3~cups, is the sauce to keep.
\end{enumerate}

\filbreak
\section*{Creamy Enchilada Sauce}

\noindent
Shred the \textbf{Monterey~Jack}; set aside in \textit{Medium Bowl~\#1}

\begin{enumerate}
    \item Set the empty saucepan over medium heat. Melt 3~Tbsp. \textbf{salted~butter}. Whisk in 3~Tbsp. \textbf{flour} and cook for \textit{1--2~minutes}, until the roux smells toasty and looks pale tan. Keep it pale. A darker roux tastes bitter here.

    \item Whisk in 1~cup \textbf{chicken~broth} a little at a time. Simmer for \textit{2--3~minutes}, whisking, until the sauce is thick, smooth, and heavy on a spoon. It will look pasty.

    \item Whisk in 2~cups of the green sauce from \textit{Large Bowl~\#1}. Simmer for \textit{3--5~minutes}, until the sauce is one color, glossy, and a finger path clears slowly. It should pour. Simmer for \textit{2~minutes} more if it looks thin, or loosen it with a splash of water if the spoon stands up in it.

    \item Take the pot off the heat. Stir a few tablespoons of the hot sauce into the \textbf{Mexican~crema}, then stir the \textbf{crema} back into the pot. Add 1~cup of the \textbf{Monterey~Jack} (\textit{Medium Bowl~\#1}) and stir until the shreds disappear and the sauce is smooth. Set the pot over low heat for \textit{1~minute} if the cheese will not melt. Keep it below a simmer. Taste. Add \textbf{kosher~salt} only if it tastes flat.
\end{enumerate}

\filbreak
\section*{Filling}

\noindent
Chop the \textbf{cilantro}; add 2~Tbsp. to \textit{Large Bowl~\#2} and set the rest aside in \textit{Small Bowl~\#3} (garnish) ---
Have the \textbf{chicken~tinga} ready

\begin{enumerate}
    \item Add 3~cups \textbf{chicken~tinga} and 1~cup of the \textbf{Monterey~Jack} (\textit{Medium Bowl~\#1}) to \textit{Large Bowl~\#2}, with the diced \textbf{white~onion}, the charred \textbf{corn}, and the 2~Tbsp. \textbf{cilantro} already there. Fold in 2--3~Tbsp. of the plain green sauce from \textit{Large Bowl~\#1}. The meat should look moist and coated. Add another spoonful of green sauce if the \textbf{tinga} looks dry. Taste, and add \textbf{kosher~salt} only if it tastes flat.
\end{enumerate}

\filbreak
\section*{Enchiladas}

\noindent
Heat the oven to \textit{375°F} ---
Crumble the \textbf{cotija} into \textit{Small Bowl~\#3}, with the remaining \textbf{cilantro}

\begin{enumerate}
    \item Spread about \nicefrac{3}{4}~cup of the creamy sauce in a 9~inch$\times$13~inch baking dish. Pour about 1~cup more into a shallow bowl for dipping. Leave the rest in the saucepan.

    \item While the oven heats, set the skillet over medium heat. Add the remaining \textbf{vegetable~oil} and heat until it shimmers. Fry the \textbf{tortillas} one at a time for \textit{5--10~seconds} per side, until they soften, blister slightly, and bend without cracking. Drain them and stack them in a towel. Continue for a few seconds per side if a tortilla still cracks. Add a spoonful of \textbf{oil} if the skillet goes dry.

    \item Dip each \textbf{tortilla} in the creamy sauce, let the excess drip off, and spoon about 3~Tbsp. of the filling (\textit{Large Bowl~\#2}) down the center. Roll it closed and set it seam-side down in the dish. Pack the rolls so they hold each other shut. Extra filling keeps for another meal.

    \item Pour the remaining creamy sauce over the enchiladas, including any sauce left in the dipping bowl. Scatter about \nicefrac{1}{2}~cup of the \textbf{Monterey~Jack} (\textit{Medium Bowl~\#1}) over the top. Cover with foil.

    \item Bake at \textit{375°F}, covered, for \textit{15~minutes} when the filling went in warm, or \textit{25~minutes} when it was cold from the refrigerator.

    \item Uncover and bake until the \textbf{Monterey~Jack} is melted, the sauce bubbles at the edges, and the center of the dish reads \textit{165°F}, about \textit{8--12~minutes} more. Continue in \textit{3~minute} increments until the center is hot and the edges are bubbling.

    \item Rest for \textit{5~minutes}. Scatter the \textbf{cotija} and the remaining \textbf{cilantro} (\textit{Small Bowl~\#3}) over the top. The enchiladas are at their best in the next half hour. The plain green sauce in \textit{Large Bowl~\#1} is the batch to keep.
\end{enumerate}

\newpage

{\small
\setlength{\columnsep}{20pt}
\setlength{\multicolsep}{6pt}
\begin{multicols}{2}
\setlength{\parindent}{0pt}
\setlength{\parskip}{4pt}

\section*{Equipment Required}
\begin{itemize}
    \item Air fryer
    \item Blender
    \item Large saucepan, 3- to 4-quart
    \item 12-inch skillet
    \item 9~inch$\times$13~inch baking dish
    \item Shallow bowl or pie plate, for dipping
    \item Colander
    \item Whisk
    \item Tongs
    \item Wooden spoon or spatula
    \item Cutting board
    \item Chef's knife
    \item Measuring cups and spoons
    \item Can opener
    \item Small prep bowls (3)
    \item Medium prep bowl (1)
    \item Large prep bowls (2)
    \item Instant-read thermometer
    \item Aluminum foil
\end{itemize}

\section*{Hints and Notes}

\subsection*{Yield}
\begin{itemize}
    \item Makes 12 enchiladas in a 9~inch$\times$13~inch dish. Serves 4--6.
    \item The green sauce yields about 5~cups after it coats a spoon. The enchiladas use 2~cups. About 3~cups remain for later. A looser simmer leaves more. A longer simmer leaves less.
\end{itemize}

\subsection*{Mise en Place}
\begin{itemize}
    \item Small Bowl \#1 --- sauce spices: 1\nicefrac{1}{2}~tsp. \textbf{cumin}, 1~tsp. \textbf{Mexican~oregano}
    \item Small Bowl \#2 --- juice of 1 \textbf{lime}
    \item Small Bowl \#3 --- garnish: remaining \textbf{cilantro}, then the crumbled \textbf{cotija} (3~oz.)
    \item Medium Bowl \#1 --- shredded \textbf{Monterey~Jack} (10~oz., about 2\nicefrac{1}{2}~cups)
    \item Large Bowl \#1 --- finished plain green sauce; 2~cups go into the creamy sauce, and the rest stays here
    \item Large Bowl \#2 --- filling: \nicefrac{1}{2}~cup diced \textbf{white~onion}, then the charred \textbf{corn}, then \textbf{chicken~tinga}, 1~cup \textbf{Monterey~Jack}, and 2~Tbsp. \textbf{cilantro}
    \item Blender --- drained \textbf{tomatillos}, most of the \textbf{onion}, and the peeled \textbf{garlic}, then the seeded chiles
    \item Prep the green sauce while the chiles air-fry. Char the \textbf{corn} while the green sauce simmers. Fry the \textbf{tortillas} while the oven heats.
\end{itemize}

\subsection*{Ingredient Tips}
\begin{itemize}
    \item \textbf{Anaheim~chiles} and \textbf{poblano~chiles}: weigh them. 18~oz. of Anaheim to 6~oz. of poblano is 3:1. That is about 5--6 Anaheim and 1--2 poblano. Anaheim carries the green-chile flavor. Poblano adds body and sweetness.
    \item \textbf{Jalapeño}: one pepper, seeds scraped out after cooking. It warms the sauce. The chipotle in the \textbf{chicken~tinga} is the main heat.
    \item \textbf{Tomatillos}: one 28~oz. can, drained, liquid discarded. They go into the blender as they are. The simmer is what concentrates them.
    \item \textbf{White~onions}: two medium onions. \nicefrac{1}{2}~cup stays raw in the filling. The rest is blended into the sauce. A medium onion is about 1--1\nicefrac{1}{2}~cups diced, so the sauce gets a generous amount.
    \item \textbf{Chicken~tinga}: use the tinga in this book. Measure 3~cups of the sauced, shredded meat. One batch is more than this pan needs.
    \item \textbf{Salted~butter}: it seasons the creamy sauce. Taste before adding more salt.
    \item \textbf{Chicken~broth}: low-sodium. The green sauce, the butter, and the cheese already carry salt.
    \item \textbf{Mexican~crema}: sour cream works in its place. Temper either one so it stays smooth.
    \item \textbf{Monterey~Jack}: 10~oz. is about 2\nicefrac{1}{2}~cups shredded. 1~cup melts into the creamy sauce, 1~cup goes into the filling, and about \nicefrac{1}{2}~cup melts on top in the oven.
    \item \textbf{Cotija}: 3~oz., crumbled. It stays salty and dry. Scatter it on after the rest so it does not melt into the \textbf{Jack}.
    \item \textbf{Corn~tortillas}: six-inch corn tortillas. Ones a day old fry and roll with less breaking.
    \item \textbf{Corn~kernels}: fresh or frozen. Pat them dry so they brown after the moisture cooks off.
\end{itemize}

\subsection*{Preparation Tips}
\begin{itemize}
    \item Keep the chiles in one layer. Overlap softens them with steam, and the skins stay pale.
    \item The charred skins stay on. Blend until the flakes are gone. Those skins are the roast in this sauce.
    \item The green puree is thick and will scorch if it sits. Scrape the bottom before you leave it to char the \textbf{corn}, then stir it again as soon as the \textbf{corn} is done.
    \item After the broth goes into the roux, the sauce looks heavy. The 2~cups of green sauce are what bring it back to a pourable coat.
    \item Stir the \textbf{crema} in off the heat. A boil can make it look grainy. Low heat melts the \textbf{Jack} without splitting the sauce.
    \item Fry the \textbf{tortillas} only until they bend. A crisp tortilla cracks when it is rolled. The measured \textbf{oil} covers the chiles, the simmer, and the skillet. Keep a little more nearby if the skillet goes dry.
    \item Fill with about 3~Tbsp. A packed tortilla splits. Extra filling keeps for another meal.
    \item Cold \textbf{tinga} needs the longer covered bake so the center reaches \textit{165°F}. Warm \textbf{tinga} is ready on the shorter time.
    \item \textbf{Cotija} goes on after the \textit{5~minute} rest. Under the foil it softens into the melted \textbf{Jack}.
\end{itemize}

\subsection*{Make Ahead \& Storage}
\begin{itemize}
    \item The plain green sauce is the part to save. Cool it, then refrigerate up to \textit{4~days} or freeze up to \textit{3~months}. Thaw overnight in the refrigerator.
    \item Reheat the plain sauce over low heat. Add a splash of water if it has tightened.
    \item Make the green sauce ahead, then build the creamy sauce on the day you bake. The creamy sauce is best the day it is made.
    \item Baked enchiladas keep up to \textit{4~days} in the refrigerator. Reheat, covered, at \textit{350°F} until the center reads \textit{165°F}.
    \item Add the \textbf{cilantro} and \textbf{cotija} after reheating.
    \item The enchiladas are at their best in the half hour after the \textit{5~minute} rest.
\end{itemize}

\subsection*{Serving Suggestions}
\begin{itemize}
    \item Serve with Mexican rice and beans, diced \textbf{white~onion}, and a squeeze of \textbf{lime}.
    \item Keep the plain green sauce for eggs, pork, chicken, rice, or another pan of enchiladas.
    \item A Mexican lager or a hibiscus agua fresca stands up to the chile and the crema.
\end{itemize}

\end{multicols}
}

\end{document}
